\section{Complete public transcripts}
\label{sec:transcripts}

To apply the four-effect inequality to a protocol, we must obtain
$X_i(t)\otimes Y_j(t)$ at each public transcript $t$, with the same two
local pairs for all four bit values.  These matrices describe the effects
of the entire protocol on its input.  They are not the effects of just its
last measurements.  We first derive this form for finite classical records,
then prove its measurable counterpart for the admitted protocols below.
The transcript always retains every public message, public random draw,
and public stopping decision.

\subsection{The finite-record mechanism}

Let $R$ be the input state on finite-dimensional spaces $A\otimes B$.
In fixed product bases, Eve's canonical purification is
\begin{equation}\label{eq:canonical-purification}
 |\Psi_R\rangle=\sum_h\sqrt R\,|h\rangle_{AB}|h\rangle_E,
 \qquad E\cong A\otimes B.
\end{equation}
For the moment suppose there are finitely many local-instrument and public
steps, including the final readouts, every classical record is finite, and
each party produces its prescribed bit almost surely on $R$.  Use the
following resource rules: the recorded initial public variable
$C_0$ is independent of the input and reference; conditional on $C_0$, that
input-reference state and the two private auxiliary resources form a product.
The initial public law and local preparation kernels are fixed when the
input varies.  All later public-source draws use fixed input-independent
kernels of the preceding public record.  Thus the input is the only shared private resource.

Fix a complete public transcript $t$.  In Alice's laboratory, compose the
instrument branches consistent with $t$ and final bit $i$, retaining private
memory until it is no longer used and summing private outcomes locally.
This gives a CP map $\mathcal A_i^t$ on her input.  Bob obtains a CP map
$\mathcal B_j^t$ in the same way.  At this fixed public record each local
choice can depend on the preceding public messages but not on the other
party's unannounced private outcomes.  The compatible private histories
therefore sum independently within the two laboratories.  Public-source
weights are scalars fixed by $t$; absorbing their product into either local
map, the joint branch is $\mathcal A_i^t\otimes\mathcal B_j^t$.

Set $X_i(t)=(\mathcal A_i^t)^\dagger(I)$ and
$Y_j(t)=(\mathcal B_j^t)^\dagger(I)$.  These matrices are positive and the
branch effect is $D_{ij}(t)=X_i(t)\otimes Y_j(t)$.  Tracing the final
laboratory memories out of that branch leaves Eve with
\[
 \tau_{ij}(t)=\bigl(\sqrt R\,D_{ij}(t)\sqrt R\bigr)^\T.
\]
Indeed, the $(h,k)$ matrix entry on Eve's system is
$\operatorname{tr}[D_{ij}(t)\sqrt R|h\rangle\langle k|\sqrt R]
=\langle k|\sqrt R D_{ij}(t)\sqrt R|h\rangle$.
The same $X_0,X_1$ and $Y_0,Y_1$ occur in all four branches.  Private
instrument labels have been summed, not added to Eve's transcript.
No pointwise normalization is required at a fixed $t$; normalization holds
after summing all final bits and all transcripts.

For a continuous transcript, a singleton may have probability zero, and
with countably many steps an output may depend on the complete local
history.  The general statement will replace the preceding matrices by
densities with respect to a single transcript measure.  Its proof runs the
protocol on a product of two local probes, each with an untouched reference.
On that product input the local histories are independent conditional on
the complete public record.  Reference measurements recover common product
matrix densities from this scalar independence, and a fixed filter on the
references gives the blocks for $R$.

\subsection{The admitted protocol model}

The instrument formalism originates with Davies and Lewis
\cite{DaviesLewis1970}.  Chitambar, Leung, Man\v{c}inska, Ozols, and Winter
\cite{CLMOW2014} distinguish finite-round LOCC, infinite linked protocols,
and topological closure.  We specify the measurable model used here so that
both its complete-history law and any terminal readout are part of the
protocol data.  Ordinary finite protocols obtain these laws by composition;
a countable schedule requires the compatibility conditions below.

\begin{definition}[Admitted local protocols]\label{def:protocol}
Let $A$ and $B$ be finite-dimensional complex input spaces.  An admitted
protocol consists of local quantum instruments, local memories and
randomness, and authenticated two-way public classical communication.
Each party may process its whole input space jointly.  Its data and their
joint laws satisfy the following requirements.
\begin{enumerate}
\item\label{protocol:initial} \textit{Initial resources.}
Before interaction,
there is a recorded standard Borel public variable $C_0$, possibly constant,
whose law is fixed once the protocol is chosen.  It is independent of the
supplied input together with any purifying or untouched reference.
Conditional on $C_0$, the joint state or law of that input and reference
and the two initial private auxiliary resources is the product of the
input-reference state and one resource for each laboratory.  These local
resources include private ancillas and randomness and may depend on $C_0$.
The public law and these two preparation kernels are the same whenever
this protocol is run on another input.  All shared initial randomness is
included in $C_0$.
An initialization acting on a party's input is treated as a first local
operation.  Apart from this recorded public variable and the supplied
input, there is no shared resource and no preshared secret key.

\item\label{protocol:instruments} \textit{Local instruments.}
An instrument assigns completely positive operations to measurable
outcome sets, countably additively on input states.  It has measurable
outcome laws and measurable conditional local updates, and its total
operation is trace preserving.  Its classical outcome space is standard
Borel.  Local instruments may depend measurably on the local private
history and the preceding public record.

\item\label{protocol:history} \textit{Local aggregation and chronological history.}
A local subroutine completed before a later step may be aggregated into
one instrument when that aggregate is supplied as a local instrument
reproducing its retained output and reference law on arbitrary inputs with
untouched finite-dimensional references.  It includes every private record
later used and its continuing memory.  Its total is normalized, including
a nonreturn or failure tag when needed.
All public communication remains in the public schedule.  After any such
aggregation and before a separate terminal readout on a complete infinite
history, the genuine classical history, including private instrument
outcomes and public records, forms a finite sequence or a sequence indexed
by $\mathbb N$, in fixed standard Borel coordinate spaces.  Instrument-choice
tags and fixed padding for stopped histories are included in these spaces.
A local choice uses only that laboratory's preceding private history and
the preceding public record; any choice affecting the other laboratory is
communicated through a public step.

\item\label{protocol:public} \textit{Public steps and the complete record.}
The public discussion is correspondingly finite or indexed by $\mathbb N$.
Each message is produced by its sender's local process from that sender's
private history and the preceding public record.  For a later public-source
draw after the public record $T_m=t$, the recorded outcome is drawn from a
measurable probability kernel $\kappa_m(t,dc)$ depending only on that record.
On every unnormalized joint instrument law, this public update multiplies
by the scalar kernel and acts as the identity on the quantum and reference
registers.  The kernel at a given record is the same for every initial input.
A party-generated announced choice, deterministic or randomized, is that
party's message.  Public choices of speaker and public stopping decisions
use these same two rules; a choice measurable from the preceding public
record is a Dirac public kernel.  Simultaneous messages may be ordered while
retaining their dependence on the earlier record.  Every publicly observed
control or other datum is generated and recorded by one of these steps.
The standard Borel record $T$ is the encoded sequence of $C_0$ and all
these public messages, draws, speaker choices, and stopping decisions.
Private outcomes and quantum
memories remain in their own laboratories.

\item\label{protocol:reference} \textit{Reference-compatible prefix laws.}
The specified steps are one family of instruments on arbitrary positive
trace-class inputs, with their amplifications by any untouched
finite-dimensional reference.  Their unnormalized branch laws give
countably additive positive matrix-valued measures on that reference for
finite prefixes.  Integrating the next record recovers the preceding
reference measure after the current laboratory memories are traced out.

\item\label{protocol:terminal} \textit{Terminal outputs and local quantum readouts.}
A terminal classical output may be a measurable function of that
laboratory's instrument history and the received complete transcript $T$.
A decision to invoke a terminal quantum readout must be measurable in
the same local records.

For a terminal quantum readout, the protocol supplies both its incoming
local memory and the instrument applied to it.  The memory is supplied by
one fixed measurable local CP process on those histories, defined on
arbitrary inputs with any untouched finite-dimensional reference.  The
readout is a specified measurable local instrument on that memory,
selected using only that laboratory's private history and $T$.
The two laboratories refine the reference measure on the complete
history by these local operations, acting as the identity on the
reference.

Together with the preceding schedule, these refinements give one
countably additive, unnormalized instrument law on the complete history
and the terminal outcomes, for all initial inputs and finite untouched
references.  Include a tag for a missing or noninvoked readout among
each laboratory's outcomes.  Summing either laboratory's terminal
outcomes removes that local refinement.  Summing both and tracing the
output memories recovers the reference measure on the complete history
with its stated finite-prefix marginals.

After any terminal outcomes have been included, each prescribed output
status is a measurable local value in the finite output alphabet
together with a symbol $\perp$ for nonproduction.  An output already
produced is retained.  The incoming-memory and compatibility requirements
concern only memory used by a terminal readout; unused memory need not
converge.  An unspecified evolution through infinitely many quantum
operations does not by itself define a terminal readout.

\item\label{protocol:completion} \textit{Completion on a specified input.}
A \emph{completed bit-output protocol on a specified input} has a bit
in each laboratory almost surely on that input, with normalized joint
output.  Almost-sure production is required only on this actual input.
More generally a completed key-output protocol produces an output in a
specified finite classical alphabet at each party almost surely on the
actual input, again with normalized joint output.  Trace-norm limits of
completed actual outputs, with full recoverable public information, will be
treated separately in Section~\ref{sec:secret-key}.
\end{enumerate}
\end{definition}

For a finite protocol whose instruments and public steps satisfy
clauses~\ref{protocol:initial}--\ref{protocol:public},
composition of those instruments supplies the prefix laws
and any finite terminal refinement in
clauses~\ref{protocol:reference}--\ref{protocol:terminal}.  The separate
complete-history and terminal-process requirements matter when a countable
schedule or an aggregated subroutine is used.
There is no uniform bound on the amount of communication, the size of a
standard Borel message, or the number of rounds within this scope.  In
particular, a terminal classical bit can depend on an entire infinite
record.

\subsection{Common effects at a complete transcript}

\begin{lemma}[Common local factors for a complete transcript]
\label{lem:transcript}
Let $R$ be a density operator on finite-dimensional complex spaces
$A\otimes B$, and let an admitted protocol produce bits
$i,j\in\{0,1\}$ almost surely on $R$.  Give Eve the canonical
purification~\eqref{eq:canonical-purification},
and let her retain $E$ and the complete public transcript $T$ with value
$t$ in its standard Borel space $\mathsf T$.  There exist a finite
positive measure $\mu$ on $\mathsf T$ and measurable positive matrices
$X_0(t),X_1(t)$ on $A$ and $Y_0(t),Y_1(t)$ on $B$ such that the positive
matrix-valued measure of Eve in bit branch $(i,j)$ has the density
\begin{equation}\label{eq:eve-blocks}
 \tau_{ij}(t)=
 \bigl(\sqrt R\,[X_i(t)\otimes Y_j(t)]\sqrt R\bigr)^\T
 \quad\text{with respect to }\mu.
\end{equation}
The same $X_i(t)$ and $Y_j(t)$ work for all four branches outside one
$\mu$-null set, and
\[
 \sum_{i,j=0}^1\int_{\mathsf T}\tr\tau_{ij}(t)\,d\mu(t)=1.
\]
The transpose in \eqref{eq:eve-blocks} is in the product basis
identifying $E$ with $A\otimes B$.  No invertibility of $R$ is assumed.
\end{lemma}

\begin{proof}
Write $a=\dim A$ and $b=\dim B$, and attach untouched finite references
$A_0\cong A$ and $B_0\cong B$.  Run the protocol on the product of the
normalized local probes
\begin{equation}\label{eq:local-probes}
 |\Phi_a\rangle_{AA_0}\otimes|\Phi_b\rangle_{BB_0},
 \qquad
 |\Phi_a\rangle=a^{-1/2}\sum_{r=1}^{a}|r,r\rangle,
 \qquad
 |\Phi_b\rangle=b^{-1/2}\sum_{s=1}^{b}|s,s\rangle.
\end{equation}
The protocol need not produce both bits on this auxiliary input.  For
this experiment assign bit zero in a laboratory if its prescribed bit
is never produced, and retain a bit that has already been produced.
The assignment is a measurable local rule.  It does not replace the
other laboratory's output when only one output is missing.  We will
remove the contribution of these auxiliary defaults after recovering
the actual input.

The proof has three steps.  First, clauses~\ref{protocol:initial}--\ref{protocol:public}
give conditional independence of the two local sampling tapes on the
product probes, given the complete public transcript.  Next,
clauses~\ref{protocol:reference}--\ref{protocol:terminal} give positive
reference measures; informationally complete tests turn the scalar
independence into product matrix densities.  Finally, a filter on the
untouched references prepares \eqref{eq:canonical-purification}
and transfers those densities to Eve's actual blocks.

\paragraph{Conditional independence on the product probes.}
The scalar factorization has the finite private-coin rectangle property
as a classical antecedent \cite[Lemma~6.7]{BarYossefEtAl2004}.
We prove the complete-history version needed here before recovering its
matrix-valued factors.

Choose a finite informationally complete POVM on each reference.  One
explicit construction in dimension $d$ starts with the projections
onto the basis vectors and onto
\[
 (|r\rangle+|s\rangle)/\sqrt2,
 \qquad
 (|r\rangle+\mathrm i|s\rangle)/\sqrt2
 \quad (r<s).
\]
These projections span the real vector space of Hermitian matrices.  If
they are $P_1,\ldots,P_N$, the $2N$ effects
$P_k/N$ and $(I-P_k)/N$ form a POVM with the same spanning property.
Call the two resulting POVMs $(E_r)_r$ and $(G_s)_s$.  Because the
references are untouched, these tests commute with all laboratory
operations.  Their joint statistics with the transcript and final bits
can therefore be computed with the tests performed at the start.

In this tested product-probe experiment, sample $C_0$ with the same fixed
law as in the actual run.  Conditional on $C_0$, all local outcomes can be
sampled using independent local random tapes $U$ for Alice and $V$ for Bob.
The tapes may be taken as countable products of uniform variables, with
coordinates reserved for the reference tests and any terminal readouts.
They are chosen before the sampling; the local outcome histories are then
functions of these tapes and the public record.
At each local instrument, sample
its conditional outcome law using a fresh uniform variable on that
laboratory's tape and apply the stipulated conditional state update.
The tested input and private ancillas are a product across laboratories
conditional on $C_0$, and each local update changes only its own local
state.  At a public-source step, sample $\kappa_m(T_m,dc)$ using a fresh
public uniform variable independent of the private tapes and all prior
sampling variables; only the recorded
draw enters the transcript.  This step leaves the quantum registers
unchanged.  Induction over the specified finite prefixes therefore shows
that this sampling has the genuine joint finite-prefix laws.
For a standard Borel outcome space a probability kernel has such a
jointly measurable sampler: one may code the space by a Borel subset
of $[0,1]$ and use the measurable kernel-sampling theorem
\cite[Appendix~A1, Theorem~A1.2, and Lemma~3.22]{Kallenberg2002}.
The initial reference tests and each specified local terminal quantum
readout are sampled in the same way.  Their stated local compatible
refinements give the genuine tested laws on the complete histories.
A terminal classical readout is
already a measurable function of its records.  Each defaulted bit and
each reference-test result is consequently a measurable function of
that laboratory's tape and the complete public transcript.  This
sampling describes the auxiliary product input; the actual input $R$
need not have a product simulation.

Let $T_m$ consist of $C_0$ and the first $m$ subsequent public records,
and let $\mathcal F_m=\sigma(T_m)$.  At a step sent by Alice the new
record is measurable with respect to
$\sigma(U)\vee\mathcal F_m$; at a step sent by Bob the analogous
statement holds with $V$.  At a public-source step with new draw $C$,
the sampling construction gives
\[
 \mathbb P(C\in dc\mid\sigma(U,V)\vee\mathcal F_m)=\kappa_m(T_m,dc).
\]
We show by induction that $U$ and $V$ are conditionally independent
given $\mathcal F_m$.  The assertion at $m=0$ follows from the product
initial resources and independent private tapes conditional on $C_0$.

Suppose it holds at $m$ and Alice sends the next record.  For bounded
measurable functions $u$ and $v$, conditional independence gives
\[
 \mathbb E[v(V)\mid\sigma(U)\vee\mathcal F_m]
 =\mathbb E[v(V)\mid\mathcal F_m].
\]
As $\mathcal F_{m+1}\subseteq\sigma(U)\vee\mathcal F_m$, conditioning
this equality on $\mathcal F_{m+1}$ shows that
\[
 \mathbb E[v(V)\mid\mathcal F_{m+1}]
 =\mathbb E[v(V)\mid\mathcal F_m].
\]
Using the same equality inside the conditional expectation of
$u(U)v(V)$ gives
\begin{equation}\label{eq:tape-factorization}
 \mathbb E[u(U)v(V)\mid\mathcal F_{m+1}]
 =\mathbb E[u(U)\mid\mathcal F_{m+1}]
  \mathbb E[v(V)\mid\mathcal F_{m+1}].
\end{equation}
The Bob step is symmetric.

For a public-source step, let $\varphi$ be a bounded measurable function
of its new draw $C$, and write
$(\kappa_m\varphi)(T_m)=\int\varphi(c)\,\kappa_m(T_m,dc)$.  The kernel
condition and the induction hypothesis give
\[
 \mathbb E[u(U)v(V)\varphi(C)\mid\mathcal F_m]
 =(\kappa_m\varphi)(T_m)
   \mathbb E[u(U)\mid\mathcal F_m]\,
   \mathbb E[v(V)\mid\mathcal F_m].
\]
Testing against bounded $\mathcal F_m$-measurable factors times
$\varphi(C)$ and applying the monotone class theorem identifies conditional expectations on
$\mathcal F_{m+1}=\mathcal F_m\vee\sigma(C)$.  Taking $v=1$ or $u=1$
shows that the respective conditional marginal is unchanged; taking
both functions gives their product.  Thus the same factorization
\eqref{eq:tape-factorization} holds after the public-source step.
The rules for speaker and stopping records are among these sender and
public-source steps.  If two messages were simultaneous, ordering them
preserves the required measurability for the second, because it still
depends only on its original earlier record.  Fixed padding after a
public stop also preserves the claim.

Set $\mathcal F_\infty=\sigma(\bigcup_m\mathcal F_m)=\sigma(T)$, the sigma
field of the actual complete transcript.  The auxiliary sampling uniforms
for later public draws are not adjoined.  If $Z$ is bounded, then
\[
 \mathbb E[Z\mid\mathcal F_m]\longrightarrow
 \mathbb E[Z\mid\mathcal F_\infty]\quad\text{in }L^2.
\]
Indeed, these are the orthogonal projections of $Z$ onto increasing
closed subspaces of $L^2$.  Their limit is the projection onto the
closure of the union.  Indicators of finite-prefix events belong to
that union, and the monotone class theorem shows that their span is
dense in $L^2(\mathcal F_\infty)$.  The limiting projection is therefore
the indicated conditional expectation.  For bounded $u,v$, the two
conditional expectations on the right of \eqref{eq:tape-factorization}
converge in $L^2$ and remain bounded.  Cauchy--Schwarz applied to the
difference of their products gives convergence in $L^1$.  Passing to
the limit in \eqref{eq:tape-factorization} proves
\[
 \mathbb E[u(U)v(V)\mid\mathcal F_\infty]
 =\mathbb E[u(U)\mid\mathcal F_\infty]
  \mathbb E[v(V)\mid\mathcal F_\infty].
\]

The tape spaces and $\mathsf T$ are standard Borel.  They admit regular
conditional probability kernels and conditional integration of
functions depending on $t$
\cite[Theorems~6.3--6.4]{Kallenberg2002}.  Apply the last factorization
to indicators from countable generating algebras for the two tape spaces,
including the whole spaces, and discard the union of the resulting null
sets.  Their product rectangles form a $\pi$-system generating the product
sigma field.  Uniqueness of probability measures on that $\pi$-system then
yields, on a single set of full transcript measure,
\begin{equation}\label{eq:full-tape-law}
 \mathsf P_{U,V\mid T=t}
 =\mathsf P_{U\mid T=t}\otimes\mathsf P_{V\mid T=t}.
\end{equation}
In particular this law factors the conditional expectations of any
bounded measurable function of $(U,t)$ and any bounded measurable
function of $(V,t)$.  Thus the pair consisting of Alice's defaulted
bit and her reference-test result is conditionally independent of the
corresponding pair for Bob, given $t$.

\paragraph{Positive reference measures and their densities.}
Return to the experiment in which the probes are retained and the
reference tests are not performed.  We first construct its reference
matrix measure on the space of all genuine classical instrument
histories, including private histories.  The fixed initial law of $C_0$
and the specified finite-prefix instruments give a positive reference
matrix measure at every prefix.  Summing the next local outcome recovers
the preceding reference measure because the total instrument is trace
preserving.  Integrating a public-source outcome does the same because
its scalar kernel has total mass one and acts as the identity on the
quantum registers.  These measures are therefore consistent on finite
cylinders.

Here is the countable extension in finite-dimensional terms.  For a
positive reference test, its scalar finite-prefix laws are consistent
finite measures; after normalization, or by adjoining the
complementary test outcome, the countable product extension theorem
applies \cite[Theorem~6.14]{Kallenberg2002}.  It supplies a finite
scalar measure on the countable standard Borel history space.  Recover
matrix entries from finitely many rank-one reference tests by
polarization.  For an arbitrary reference vector $z$, extend also its
positive rank-one test.  On finite cylinders the resulting scalar
measure equals the quadratic form $z^*W(\,\cdot\,)z$ of the recovered
matrix measure.  Uniqueness of finite measures on cylinders extends
this equality to the entire history sigma field.  Thus
$z^*W(S)z\ge0$ for every measurable $S$ and every $z$, so $W$ is a
positive matrix-valued measure.  This argument constructs a measure
for the untouched finite references; it asks for no limiting state of
unused quantum memory.

Measurable terminal classical readouts and the local default rule now
give measurable events in the history space.  If a terminal quantum
readout is used, the incoming local process stipulated in
Definition~\ref{def:protocol}, on those same instrument histories and
with an untouched finite reference, together with its specified local
instrument gives the bit-conditioned positive reference measures.
Each instrument is selected from its own laboratory's history and $T$;
their joint refinement is obtained by the two local operations.  Include
the missing or noninvoked tag among the outcomes, so summing both local
outcomes recovers the reference measure on the complete history.  This
terminal operation is a separate extension of that history; its outcome is
included in the enlarged history before the classical readout.  This is where
clauses~\ref{protocol:reference}--\ref{protocol:terminal} supply the needed
locality, prefix compatibility, and reference stability.  Write
$\widehat W$ for the resulting positive reference measure on the enlarged
history space.  Its marginal on the preterminal history is $W$; when no
readout is used, identify $\widehat W$ with $W$ using trivial terminal tags.
The construction uses the genuine
instrument records; the tapes above were only a sampling device for
their tested scalar laws.

Let $W_{ij}(S)$ be the restriction and pushforward of $\widehat W$ to
the defaulted bit pair $(i,j)$ and transcript in a measurable set
$S\subseteq\mathsf T$.  Define the probe transcript law
\[
 \mu(S)=\sum_{i,j=0}^1\tr W_{ij}(S).
\]
It is a probability measure: the probe is normalized, the instrument
totals preserve trace, and the defaults assign a pair on every history.
For each $i,j$, positivity and $\tr W_{ij}(S)\le\mu(S)$ imply absolute
continuity of every matrix entry with respect to $\mu$.  Hence there
are entrywise Radon--Nikodym densities $W_{ij}(t)$.  Only the dominated
finite-measure case is needed here.  For completeness, if
$0\le\nu\le c\mu$, the functional $g\mapsto\int g\,d\nu$ on simple
functions is bounded in $L^2(\mu)$ by Cauchy--Schwarz.  Its Hilbert
space representing vector is a density, and indicator tests show that
the density lies between $0$ and $c$.  Apply this fact to the finitely
many polarization tests for the matrix entries.

Testing the densities on a countable dense set of reference vectors
shows that $W_{ij}(t)\succeq0$ outside one null set.  The equality
defining $\mu$ also gives
\[
 \sum_{i,j}\tr W_{ij}(t)=1\quad\text{for $\mu$-almost every }t.
\]
Put
\[
 U_i(t)=\sum_j\tr_{B_0}W_{ij}(t),
 \qquad
 V_j(t)=\sum_i\tr_{A_0}W_{ij}(t).
\]
Both local pairs are positive, and
$\sum_i\tr U_i(t)=\sum_j\tr V_j(t)=1$ almost everywhere.

The conditional scalar statistics of the reference tests are the
traces against these densities.  Their local marginals are
$\tr[E_rU_i(t)]$ and $\tr[G_sV_j(t)]$, because the other POVM sums to
the identity.  Conditional independence from
\eqref{eq:full-tape-law} therefore gives
\[
 \tr[(E_r\otimes G_s)W_{ij}(t)]
 =\tr[E_rU_i(t)]\,\tr[G_sV_j(t)]
\]
for every $i,j,r,s$, outside one null set.  The tensor products
$E_r\otimes G_s$ span the Hermitian matrices on $A_0\otimes B_0$.
Equality against that spanning family proves
\begin{equation}\label{eq:probe-product}
 W_{ij}(t)=U_i(t)\otimes V_j(t)
 \quad\text{for all }i,j\text{ and $\mu$-almost every }t.
\end{equation}
We have obtained the same two local matrix pairs for all four branches,
without conditioning on a positive-probability singleton transcript
and without making private instrument labels public.

\paragraph{Recovering the actual purification.}
Identify $A_0\otimes B_0$ with $E$ in the product basis and define
\[
 K=\sqrt{ab}\,(\sqrt R)^\T.
\]
The elementary identity
$(I\otimes S^\T)\sum_h|h,h\rangle=(S\otimes I)\sum_h|h,h\rangle$
shows that applying $I_{AB}\otimes K$ to the product probes
\eqref{eq:local-probes} gives exactly $|\Psi_R\rangle$.
This joint filter on the references is an algebraic preparation of
the input purification; the laboratories are not required to implement
it.  It need not be trace preserving or a contraction.

The same independent law of $C_0$ and the same $C_0$-conditioned local
auxiliary resources are used on the probes and on the actual input.
Consequently the filter also prepares the full initial joint law,
including $C_0$ and those resources.
On each finite cylinder of genuine histories, the reference matrix
for this filtered input is $K W(\,\cdot\,)K^*$, because the references
are untouched and the filter commutes with every local instrument.
Every public-source update also commutes with the filter: at a given
public record it multiplies both laws by the same scalar kernel and
acts as the identity on the reference.
The scalar entries on both sides are finite measures, so uniqueness
extends their equality from cylinders to the full history sigma field.
The same commutation holds for any appended terminal instrument by
clause~\ref{protocol:terminal}.  Thus the filtered law on the enlarged
history is $K\widehat W(\,\cdot\,)K^*$, and the identity holds for the
bit-conditioned measures as well.

Let $N$ be the event in this enlarged history space that at least one
prescribed local bit is missing before the default rule.  The event may
depend on private records.  On the actual input its probability is
zero by hypothesis, so positivity gives
\[
 \tr[K\widehat W(N)K^*]=0,
 \qquad K\widehat W(N)K^*=0.
\]
Every subevent of $N$ has a positive filtered matrix bounded by this
zero matrix.  It therefore also contributes zero.  Thus the local
defaults leave the actual output measure unchanged, even though they
may occur on the probes.  The argument was made on full histories
before discarding private records, so $N$ need not be public.

It follows from \eqref{eq:probe-product} that the actual Eve density
in branch $(i,j)$ is $K(U_i(t)\otimes V_j(t))K^*$.  The matrix
$(\sqrt R)^\T$ is Hermitian.  Hence direct multiplication gives
\[
 \begin{aligned}
 K(U_i\otimes V_j)K^*
 &=ab(\sqrt R)^\T(U_i\otimes V_j)(\sqrt R)^\T\\
 &=\bigl(\sqrt R\,[aU_i^\T\otimes bV_j^\T]\sqrt R\bigr)^\T.
 \end{aligned}
\]
Take $X_i(t)=aU_i(t)^\T$ and $Y_j(t)=bV_j(t)^\T$, redefining them on a
null set if necessary.  Transposition preserves positivity, so these
are measurable positive local matrices.  The filtered total trace is
one because $R$ is normalized, and \eqref{eq:eve-blocks} follows.
No inverse of $R$ appears anywhere in the construction.
\end{proof}

The canonical purification is convenient for the formula, but it does
not restrict the adversary or the ideal comparison law.  The following
observation will let the secret-bit proof use that formula for every
purifying space.

\begin{lemma}[Recovery from another purifying space]
\label{lem:purification-recovery}
Let $R$ and $|\Psi_R\rangle$ be as in Lemma~\ref{lem:transcript}, and
let $|\psi\rangle\in A\otimes B\otimes\widetilde E$ be any purification
of $R$.  There is a channel $\mathcal R$ from $\widetilde E$ to $E$
that recovers the canonical purification.  If Alice and Bob run the
same admitted protocol and Eve retains her purification and the full
transcript, $\mathcal R$ acting on the purifying register and the
identity acting on the transcript recover the canonical output of
Lemma~\ref{lem:transcript}.  Moreover this channel maps every ideal
secret-bit law with Eve in $\widetilde E T$ to an ideal secret-bit law
with Eve in $ET$.
\end{lemma}

\begin{proof}
Let $\mathcal S=\supp(R^\T)\subseteq E$.  Schmidt decomposition gives
an isometry $V:\mathcal S\to\widetilde E$ with
$|\psi\rangle=(I_{AB}\otimes V)|\Psi_R\rangle$; the canonical vector
is supported on $\mathcal S$ in its purifying factor.  Choose any
density operator $\zeta$ on $E$, and define, with the compression term
included into $E$ from $\mathcal S$,
\[
 \mathcal R(\xi)=V^*\xi V+
       \tr[(I_{\widetilde E}-VV^*)\xi]\,\zeta.
\]
Compression is CP, as is the positive trace functional followed by
preparation of $\zeta$.  Their sum preserves trace, so $\mathcal R$ is
a channel.  Its first term inverts $V$ on the actual purifying support,
and the second term vanishes there.  Local operations on $A,B$ keep
each conditional purifying state supported in $\ran V$, and public-source
steps act as the identity on that register, so recovery is
exact in every transcript and bit branch.  The same formula is valid when $\widetilde E$ is
larger than the finite purifying support, acting on trace-class inputs.

An ideal bit has the form
$\frac12\sum_{i=0}^1|i,i\rangle\langle i,i|\otimes\sigma$.
Applying $\mathcal R\otimes\mathrm{id}_T$ to its Eve factor replaces
$\sigma$ by another normalized state and leaves that form intact.
This includes ideal states with mass outside $\ran V$.  Thus
trace-norm contraction transfers any lower bound proved for the
canonical output and every canonical ideal law to the output for the
arbitrary purification and every ideal law in its larger space.
\end{proof}
